% ==============================================================================
% dxh_intro.tex — DXH workshop: opening, overview, safety
% DXH講座: オープニング・全体像・安全ルール
% ==============================================================================

% ------------------------------------------------------------------
\begin{frame}{講師紹介}
  \begin{block}{伊藤 恒平 --- 金沢工業大学 ロボティクス学科 教員}
    \small StampFly 開発者 ／ 専門はドローン制御工学
  \end{block}

  \vspace{0.4em}

  \begin{table}
    \centering\small
    \begin{tabular}{l>{\raggedright\arraybackslash}p{92mm}}
      \hline
      経歴 & 大学で航空宇宙工学を学ぶ \\
           & 陸上自衛隊に入隊、飛翔体の研究開発に従事 \\
           & 高専教員を経て、金沢工業大学 ロボティクス学科 教員 \\
      \hline
      趣味 & ロボット作り・読書・映画鑑賞 \\
      \hline
    \end{tabular}
  \end{table}
\end{frame}

% ------------------------------------------------------------------
\begin{frame}{本日の流れ}
  \begin{table}
    \centering\small
    \begin{tabular}{ll}
      \hline
      \textbf{時間} & \textbf{内容} \\
      \hline
      13:30 -- 13:40（10分） & オープニング \\
      13:40 -- 14:05（25分） & ① 操縦体験（シミュレータ→実機） \\
      14:05 -- 14:30（25分） & ② 開発環境+書き込み体験 \\
      14:30 -- 15:15（45分） & ③ プログラム書き換え・モータの制御 \\
      15:15 -- 15:30（15分） & まとめ \\
      \hline
    \end{tabular}
  \end{table}

\end{frame}

% ------------------------------------------------------------------
\begin{frame}{StampFly とは}
  \begin{columns}[c]
    \begin{column}{0.45\textwidth}
      \centering
      \includegraphics[width=\columnwidth]{Stampfly_onhand.jpg}
    \end{column}
    \begin{column}{0.52\textwidth}
      \begin{table}
        \centering\small
        \begin{tabular}{ll}
          \hline
          \textbf{項目} & \textbf{仕様} \\
          \hline
          MCU    & ESP32-S3 (M5Stamp S3) \\
          IMU    & BMI270（6軸 400\,Hz） \\
          気圧   & BMP280 \\
          ToF    & VL53L3CX \\
          質量   & 37\,g（バッテリ含む） \\
          通信   & ESP-NOW + WiFi \\
          バッテリ & LiPo 1S 3.7\,V \\
          \hline
        \end{tabular}
      \end{table}
    \end{column}
  \end{columns}

  \vspace{0.3em}

  \centering {\small StampFly は \textbf{M5Stack 社 CEO の Jimmy 氏}と講師の\textbf{共同開発}です}
\end{frame}

% ------------------------------------------------------------------
\begin{frame}{StampFly の搭載センサ・部品}
  \centering
  \begin{adjustbox}{max width=0.98\textwidth, max height=0.82\textheight, center}
    \includegraphics{stampfly_spec.png}
  \end{adjustbox}
\end{frame}

% ------------------------------------------------------------------
\begin{frame}{StampFly と Tello の違い}
  \begin{center}
    \begin{tabular}{c@{\hspace{10mm}}c}
      \includegraphics[height=28mm]{tello.jpg} &
      \includegraphics[height=28mm]{stampfly_upside.jpg} \\
      {\scriptsize Tello EDU（約87\,g・カメラ搭載）} &
      {\scriptsize StampFly（37\,g）} \\
    \end{tabular}
  \end{center}

  \vspace{0.2em}

  \begin{table}
    \centering\footnotesize
    \begin{tabular}{lp{40mm}p{48mm}}
      \hline
       & \textbf{Tello} & \textbf{StampFly} \\
      \hline
      プログラミング & 飛行コマンドを外から送る（Scratch / Python） & 機体内部の制御プログラム自体を書き換える \\
      ソフトの中身 & 非公開（変更できない） & 全ソースコード公開 \\
      センサのデータ & SDK 経由の一部のみ & 生データに直接アクセス可 \\
      \hline
    \end{tabular}
  \end{table}

  \vspace{0.2em}

  \centering {\small \textbf{今日の③「モータの制御」は、中身を書き換えられる StampFly ならではの体験}}
\end{frame}

% ------------------------------------------------------------------
\begin{frame}{StampFly の現在地と、今日の楽しみ方}
  \begin{block}{StampFly はいま、育ちざかりの教材です}
    \small 位置制御まで飛ぶようになり、教材の土台が整ってきました。\\
    ここから「授業でそのまま使えるパッケージ」へ仕上げていく段階です
  \end{block}

  \vspace{0.2em}

  \begin{itemize}\small
    \item 今日は、完成すると\textbf{見えなくなる開発の生の姿}（コマンド・ビルド・調整の途中）まで体験できます
    \item ドローンの「中身」に手を入れられるのは、いまの StampFly ならではの楽しみです
  \end{itemize}

  \vspace{0.2em}

  \begin{exampleblock}{目指しているもの}
    \small \textbf{初等教育から高等教育・社会人教育まで}を見据えた教材群の完成 --- 今日はその種まきに、皆さんをお招きしています
  \end{exampleblock}
\end{frame}

% ------------------------------------------------------------------
\begin{frame}{ドローンは「勝手に」飛んでいるのではない}
  \begin{block}{Tello はなぜ「言うことを聞く」のか}
    \small 「前に進め」で前に進み、「回れ」で回る --- それは膨大な\textbf{試行錯誤（制御工学）の賜物}です。\\
    素のドローンは、制御なしでは\textbf{1秒も浮いていられません}
  \end{block}

  \vspace{0.3em}

  \begin{itemize}\small
    \item クアッドコプタは本質的に不安定 --- 放っておくと勝手に傾いて墜落する
    \item センサで姿勢を測り、4つのモータを絶えず調整し続けて、はじめて浮いていられる
  \end{itemize}

  \vspace{0.3em}

  \begin{exampleblock}{この講座の目標}
    \small \textbf{「飛ばないドローン」を制御して飛ばせるエンジニアを育てる} --- StampFly はそのための教材
  \end{exampleblock}
\end{frame}

% ------------------------------------------------------------------
\begin{frame}{クアッドコプタはどうやって飛ぶのか}
  \begin{columns}[c]
    \begin{column}{0.47\textwidth}
      \centering
      \begin{adjustbox}{max width=\columnwidth, max height=0.72\textheight, center}
        % ==============================================================================
% dxh_flight_principle_3d.tex — DXH workshop: quadcopter flight principle (3D)
% DXH講座: クアッドコプタ飛行原理の3D図
%
% Isometric-style view of an X-quad: four rotor discs with thrust vectors,
% spin-direction arcs (CW/CCW), and the three body axes with roll/pitch/yaw
% moment arrows. Projection: forward = (-0.78, 0.59), right = (0.98, 0.20),
% up = (0, 1); rotor arm length a = 1.55.
% X型クアッドの斜視図: 4つのロータ円盤+推力ベクトル、回転方向の円弧
% (CW/CCW)、機体3軸とロール/ピッチ/ヨーのモーメント矢印。
% 投影: 前方=(-0.78,0.59), 右=(0.98,0.20), 上=(0,1)、アーム長 a=1.55。
% ==============================================================================
\documentclass[border=10pt]{standalone}
\usepackage{tikz}
\usetikzlibrary{arrows.meta}
\usepackage{luatexja}
\usepackage[match]{luatexja-fontspec}
% Cross-platform font: Hiragino (macOS) → Yu Gothic (Windows) → Noto (Linux)
\usepackage{ifthen}
\newboolean{jfontset}\setboolean{jfontset}{false}
\IfFontExistsTF{Hiragino Kaku Gothic ProN}{%
  \setmainjfont{Hiragino Kaku Gothic ProN}%
  \setsansjfont{Hiragino Kaku Gothic ProN}%
  \setboolean{jfontset}{true}}{}
\ifthenelse{\boolean{jfontset}}{}{%
  \IfFontExistsTF{Yu Gothic}{%
    \setmainjfont{Yu Gothic}%
    \setsansjfont{Yu Gothic}%
    \setboolean{jfontset}{true}}{}}
\ifthenelse{\boolean{jfontset}}{}{%
  \IfFontExistsTF{Noto Sans CJK JP}{%
    \setmainjfont{Noto Sans CJK JP}%
    \setsansjfont{Noto Sans CJK JP}}{}}

\begin{document}

% Colors defined after \begin{document} (standalone compatibility, rule T9)
% 色定義は \begin{document} の後（standalone 互換, ルール T9）
\definecolor{cthrust}{RGB}{180, 60, 40}   % thrust vectors / 推力
\definecolor{ccw}{RGB}{0, 100, 180}       % CW rotor spin / CW回転
\definecolor{cccw}{RGB}{200, 90, 40}      % CCW rotor spin / CCW回転
\definecolor{cmom}{RGB}{40, 140, 60}      % moment arrows / モーメント

\begin{tikzpicture}[
    lbl/.style={font=\small\bfseries, align=center},
    axis/.style={dashed, black!45, line width=0.8pt},
    thrust/.style={-{Stealth[length=3mm]}, cthrust, line width=1.6pt},
    mom/.style={-{Stealth[length=2.6mm]}, cmom, line width=1.4pt},
    spin/.style={-{Stealth[length=2mm]}, line width=1.1pt},
  ]

  % Rotor centers (isometric projection) / ロータ中心（斜視投影）
  % FL=(-2.17,0.31)  FR=(0.62,0.87)  RL=(-0.62,-0.87)  RR=(2.17,-0.31)

  % === Body axes (behind everything) / 機体3軸（最背面） ===
  \draw[axis] (0,0) -- (-3.05,2.3);            % forward axis / 前後軸
  \draw[axis] (0,0) -- (3.15,0.64);            % lateral axis / 左右軸
  \draw[axis] (0,0) -- (0,-2.05);              % vertical axis / 上下軸

  % === Arms / アーム ===
  \draw[line width=2.4pt, black!65] (0,0) -- (-2.17,0.31);
  \draw[line width=2.4pt, black!65] (0,0) -- (0.62,0.87);
  \draw[line width=2.4pt, black!65] (0,0) -- (-0.62,-0.87);
  \draw[line width=2.4pt, black!65] (0,0) -- (2.17,-0.31);

  % === Rotor discs / ロータ円盤 ===
  \draw[fill=black!12, black!60] (-0.62,-0.87) ellipse [x radius=0.62, y radius=0.24];
  \draw[fill=black!12, black!60] (2.17,-0.31) ellipse [x radius=0.62, y radius=0.24];
  \draw[fill=black!12, black!60] (-2.17,0.31) ellipse [x radius=0.62, y radius=0.24];
  \draw[fill=black!12, black!60] (0.62,0.87) ellipse [x radius=0.62, y radius=0.24];

  % === Center body + nose / 機体中心と機首 ===
  \draw[fill=black!30, black!60] (0,0) ellipse [x radius=0.34, y radius=0.16];
  \node[lbl, text=black!70] at (-1.35,1.42) {前};

  % === Spin direction arcs / 回転方向の円弧 ===
  % CW rotors: FL (M4), RR (M2) — blue / CWロータ: FL・RR（青）
  \draw[spin, ccw] ([shift={(-2.17,0.31)}] 210:0.8 and 0.34)
    arc[x radius=0.8, y radius=0.34, start angle=210, end angle=-30];
  \draw[spin, ccw] ([shift={(2.17,-0.31)}] 210:0.8 and 0.34)
    arc[x radius=0.8, y radius=0.34, start angle=210, end angle=-30];
  % CCW rotors: FR (M1), RL (M3) — orange / CCWロータ: FR・RL（橙）
  \draw[spin, cccw] ([shift={(0.62,0.87)}] -30:0.8 and 0.34)
    arc[x radius=0.8, y radius=0.34, start angle=-30, end angle=210];
  \draw[spin, cccw] ([shift={(-0.62,-0.87)}] -30:0.8 and 0.34)
    arc[x radius=0.8, y radius=0.34, start angle=-30, end angle=210];
  \node[font=\scriptsize\bfseries, text=ccw]  at (3.05,-0.75) {CW};
  \node[font=\scriptsize\bfseries, text=cccw] at (-1.9,-1.35) {CCW};

  % === Thrust vectors / 推力ベクトル ===
  \draw[thrust] (-2.17,0.58) -- ++(0,1.15);
  \draw[thrust] (0.62,1.14) -- ++(0,1.15);
  \draw[thrust] (-0.62,-0.60) -- ++(0,1.15);
  \draw[thrust] (2.17,-0.04) -- ++(0,1.15);
  \node[lbl, text=cthrust] at (1.35,2.35) {推力};

  % === Moment arrows around each axis / 各軸まわりのモーメント矢印 ===
  % Roll (around forward axis) / ロール（前後軸まわり）
  \begin{scope}[rotate around={143:(-2.85,2.15)}]
    \draw[mom] ([shift={(-2.85,2.15)}] 185:0.38 and 0.5)
      arc[x radius=0.38, y radius=0.5, start angle=185, end angle=-40];
  \end{scope}
  \node[lbl, text=cmom] at (-3.55,1.45) {ロール};

  % Pitch (around lateral axis) / ピッチ（左右軸まわり）
  \begin{scope}[rotate around={11:(2.95,0.60)}]
    \draw[mom] ([shift={(2.95,0.60)}] 185:0.38 and 0.5)
      arc[x radius=0.38, y radius=0.5, start angle=185, end angle=-40];
  \end{scope}
  \node[lbl, text=cmom] at (3.6,-0.05) {ピッチ};

  % Yaw (around vertical axis) / ヨー（上下軸まわり）
  \draw[mom] ([shift={(0,-1.75)}] 200:0.6 and 0.24)
    arc[x radius=0.6, y radius=0.24, start angle=200, end angle=-30];
  \node[lbl, text=cmom] at (1.05,-1.85) {ヨー};

\end{tikzpicture}

\end{document}

      \end{adjustbox}
    \end{column}
    \begin{column}{0.51\textwidth}
      \begin{table}
        \centering\footnotesize
        \setlength{\tabcolsep}{3pt}
        \begin{tabular}{l>{\raggedright\arraybackslash}p{32mm}}
          \hline
          \textbf{動き} & \textbf{作り方} \\
          \hline
          浮く・上下 & 4つの推力の合計を増やす／減らす \\
          ロール（左右に傾く） & 左右の推力の差 \\
          ピッチ（前後に傾く） & 前後の推力の差 \\
          ヨー（機首を回す） & CW 組と CCW 組の回転数の差\newline（反トルク） \\
          \hline
        \end{tabular}
      \end{table}
    \end{column}
  \end{columns}

  \vspace{0.3em}

  \centering {\small 傾けば、傾いた方向へ進む。この推力配分を\textbf{1秒に400回}自動調整し続けるのが「制御」}
\end{frame}

% ------------------------------------------------------------------
\begin{frame}{今日つかう道具一式（StampFly エコシステム）}
  \centering
  \vspace{1mm}
  \begin{adjustbox}{max width=0.96\textwidth, max height=0.82\textheight, center}
    % ==============================================================================
% dxh_toolset.tex — DXH workshop: today's toolset (teacher-friendly ecosystem)
% DXH講座: 今日つかう道具一式（高校教員向けに平易化したエコシステム図）
%
% The repo-directory ecosystem_layers diagram is developer-oriented and
% meaningless to high-school teachers; this shows the same idea as the three
% physical tools on the desk plus the freely published materials.
% リポジトリ構成ベースの ecosystem_layers は開発者向けで高校教員には伝わらない
% ため、机上の3つの道具＋無料公開教材という形で同じ趣旨を示す。
% ==============================================================================
\documentclass[border=10pt]{standalone}
\usepackage{tikz}
\usetikzlibrary{arrows.meta, positioning}
\usepackage{luatexja}
\usepackage[match]{luatexja-fontspec}
% Cross-platform font: Hiragino (macOS) → Yu Gothic (Windows) → Noto (Linux)
\usepackage{ifthen}
\newboolean{jfontset}\setboolean{jfontset}{false}
\IfFontExistsTF{Hiragino Kaku Gothic ProN}{%
  \setmainjfont{Hiragino Kaku Gothic ProN}%
  \setsansjfont{Hiragino Kaku Gothic ProN}%
  \setboolean{jfontset}{true}}{}
\ifthenelse{\boolean{jfontset}}{}{%
  \IfFontExistsTF{Yu Gothic}{%
    \setmainjfont{Yu Gothic}%
    \setsansjfont{Yu Gothic}%
    \setboolean{jfontset}{true}}{}}
\ifthenelse{\boolean{jfontset}}{}{%
  \IfFontExistsTF{Noto Sans CJK JP}{%
    \setmainjfont{Noto Sans CJK JP}%
    \setsansjfont{Noto Sans CJK JP}}{}}

\begin{document}

% Colors defined after \begin{document} (standalone compatibility, rule T9)
% 色定義は \begin{document} の後（standalone 互換, ルール T9）
\definecolor{cpc}{RGB}{0, 100, 180}
\definecolor{cfly}{RGB}{180, 60, 40}
\definecolor{cjoy}{RGB}{40, 140, 60}
\definecolor{copen}{RGB}{120, 80, 160}

\begin{tikzpicture}[
    tool/.style={
      rectangle, rounded corners=5pt, draw=#1, line width=1.5pt,
      fill=#1!10, minimum width=44mm, minimum height=20mm,
      font=\bfseries\large, text=#1, align=center
    },
    desc/.style={
      font=\small, text=black!60, align=center
    },
    arr/.style={
      -{Stealth[length=3.5mm]}, line width=1.5pt, black!60
    },
    arrlabel/.style={
      font=\small, text=black!55, align=center
    },
  ]

  % === Three tools on the desk / 机の上の3つの道具 ===
  \node[tool=cpc]  (pc)  at (-6.6, 0) {ノートPC\\[-1pt]{\normalsize プログラムを書く}};
  \node[tool=cfly] (fly) at (0, 0)    {StampFly 機体\\[-1pt]{\normalsize プログラムで飛ぶ}};
  \node[tool=cjoy] (joy) at (6.6, 0)  {コントローラ\\[-1pt]{\normalsize スティックで操縦}};

  \node[desc, below=2mm of pc]  {開発環境はインストール済み};
  \node[desc, below=2mm of fly] {センサ＋小型コンピュータ搭載};
  \node[desc, below=2mm of joy] {Atom JoyStick};

  % === Horizontal connections / 横方向の接続（直角ルーティング: 水平のみ） ===
  \draw[arr] (pc.east) -- (fly.west)
    node[arrlabel, midway, above=1mm] {USBで\\書き込み};
  \draw[arr] (joy.west) -- (fly.east)
    node[arrlabel, midway, above=1mm] {電波で\\操縦};

  % === Everything is published for free / すべて無料公開 ===
  \node[rectangle, rounded corners=5pt, draw=copen, line width=1.5pt,
        fill=copen!8, minimum width=152mm, minimum height=15mm,
        align=center, below=13mm of fly, text=copen]
    {\bfseries\large ソースコード・教材・シミュレータはすべて Web で無料公開\\[0pt]
     {\normalsize\mdseries\color{black!60} 学校でも今日と同じ環境をそのまま用意できる}};

\end{tikzpicture}

\end{document}

  \end{adjustbox}
\end{frame}

% ------------------------------------------------------------------
\begin{frame}{安全ルール：飛行の安全}
  \begin{block}{飛行前の確認}
    \begin{itemize}\small\setlength{\itemsep}{0pt}
      \item 保護メガネを全員着用
      \item 飛行エリアは区画・ネットで分離
      \item プロペラガードの状態を確認
      \item 長い髪はまとめる
      \item 機体の真上に顔を出さない
    \end{itemize}
  \end{block}

  \vspace{0.15em}

  \begin{alertblock}{緊急停止}
    \footnotesize \warn{DISARM} = モータを止める操作。飛行中は自動着陸 → もう一度で即時停止
  \end{alertblock}

  \vspace{0.1em}

  \begin{exampleblock}{バッテリー}
    \footnotesize 満充電 4.2\,V／3.4\,V で低電圧警告。膨らんだ電池は使用しない。充電中は目を離さない
  \end{exampleblock}
\end{frame}

